\chapter{Background and Related Work}
\label{ch:background}

\section{Distributed Microservice Systems}
\label{sec:bg_distributed_systems}

A distributed system is a set of autonomous computing nodes that cooperate
through message passing and appear as a coherent system to users
\cite{tanenbaum2007distributed, coulouris2011distributed}. Microservices apply
this principle at the application level by splitting a system into small
services, each owned, deployed, and scaled independently
\cite{newman2015building, dragoni2017microservices}. The same decomposition
that enables rapid deployment also creates operational complexity: service
instances are ephemeral, dependencies change frequently, and a user-visible
symptom may be several network calls away from its original cause.

Reliability in such systems is constrained by partial failure. The CAP theorem,
introduced by Brewer and formalized by Gilbert and Lynch, states that a
distributed data store cannot simultaneously guarantee consistency,
availability, and partition tolerance under network partitions
\cite{brewer2000towards, gilbert2002brewer}. Monitoring systems usually
prioritize availability and partition tolerance: local anomaly detection should
continue even when the global coordinator is temporarily unreachable. This
choice accepts short-lived model or graph staleness in exchange for continuous
visibility.
\label{subsec:bg_cap}

Observability is built from metrics, logs, and traces
\cite{kalloniatis2020microservices_observability}. Metrics describe numerical
health indicators such as CPU utilization, latency percentiles, memory pressure,
request rate, and error rate. Logs provide event-level context. Traces reconstruct
request paths across service boundaries and make dependency-aware diagnosis
possible \cite{sigelman2010dapper}. An effective anomaly detection framework
must combine these signals: metrics trigger detection, traces provide topology,
and historical behavior provides the baseline against which abnormality is
defined.

Existing observability platforms typically centralize raw telemetry for
analysis, which introduces bandwidth bottlenecks and privacy exposure at scale.
This thesis addresses that limitation through edge-local detection
(\S\ref{sec:method_lstm_ae}) and federated model training
(\S\ref{sec:fl_architecture}), so that raw data remains at the source while
model intelligence is shared.

\section{Anomaly Detection for Telemetry}
\label{sec:bg_anomaly}

Classical anomaly detection uses static limits, z-scores, rolling statistics, or
Exponentially Weighted Moving Averages (EWMA) \cite{basseville1993detection,
chandola2009anomaly}. These approaches are inexpensive and interpretable, but
they are weak under multivariate and non-stationary telemetry. A latency increase
may be harmless during a traffic burst, but suspicious when coupled with memory
growth and elevated downstream errors.

Deep unsupervised models are widely used when labeled failures are scarce
\cite{chalapathy2019deep, pang2021deep}. Autoencoders learn a compressed
representation of normal data and use reconstruction error as an anomaly score.
For a window $X \in \mathbb{R}^{T \times d}$ and reconstruction $\hat{X}$, the
mean squared reconstruction loss is:

\begin{equation}
    \mathcal{L}_{\text{MSE}} =
    \frac{1}{T d} \sum_{t=1}^{T}\sum_{j=1}^{d}
    \left(X_{tj} - \hat{X}_{tj}\right)^2 .
    \label{eq:bg_mse}
\end{equation}

Long Short-Term Memory (LSTM) networks are suitable for telemetry because they
capture temporal dependencies more effectively than vanilla recurrent networks
\cite{hochreiter1997long, malhotra2015lstm}. An LSTM-Autoencoder combines an
LSTM encoder, a compact latent representation, and an LSTM decoder. Windows that
do not match learned normal dynamics produce larger reconstruction errors, which
can be converted into anomaly decisions through a threshold.

However, most LSTM-based detectors assume centralized training or offline
threshold calibration. This thesis embeds an LSTM-AE within a federated
training loop (\S\ref{sec:fl_architecture}) and pairs it with adaptive
Q-learning thresholding (\S\ref{sec:threshold_implementation}), addressing
both limitations.

\section{Edge and Federated Learning}
\label{sec:bg_edge_fl}

Edge computing places computation near the data source, reducing latency,
bandwidth use, and raw-data exposure \cite{shi2016edge,
satyanarayanan2017emergence}. In this thesis, each monitored service is paired
with an edge agent that performs local preprocessing, LSTM-AE inference, and
local training. The central coordinator aggregates models rather than raw
telemetry.

Federated learning (FL) trains a shared model across clients without centralizing
their datasets \cite{mcmahan2017communication, kairouz2021advances}. Given
clients $k=1,\dots,K$ with local datasets $\mathcal{D}_k$, the global objective
is:

\begin{equation}
    \min_{\theta} F(\theta) =
    \sum_{k=1}^{K}
    \frac{|\mathcal{D}_k|}{\sum_{j=1}^{K}|\mathcal{D}_j|}
    F_k(\theta).
    \label{eq:bg_fl_objective}
\end{equation}

FedAvg updates the global model by averaging locally trained client weights:

\begin{equation}
    \theta_{t+1} =
    \sum_{k \in \mathcal{P}_t}
    \frac{n_k}{\sum_{j \in \mathcal{P}_t} n_j}
    \theta_t^{(k)},
    \label{eq:bg_fedavg}
\end{equation}

where $\mathcal{P}_t$ is the set of clients that participate in round $t$. This
partial-participation formulation is important in edge deployments, where slow
nodes, late joiners, and temporary disconnections are common.

Existing FL deployments typically assume reliable network connectivity and
ignore the communication cost of model updates. This thesis addresses both
limitations through adaptive compression (\S\ref{sec:compression}) and
pull-based model distribution (\S\ref{sec:fl_architecture}).

\section{Differential Privacy}
\label{sec:differential_privacy}

Federated learning reduces raw-data movement, but model updates can still leak
information through membership inference or model inversion attacks
\cite{shokri2017membership, fredrikson2015model}. Differential privacy (DP)
limits the influence of any single training sample on the released update
\cite{dwork2006differential, dwork2014algorithmic}. A randomized mechanism
$\mathcal{M}$ satisfies $(\varepsilon,\delta)$-DP if, for neighboring datasets
$D$ and $D'$ differing in one sample and all output sets $S$:

\begin{equation}
    \Pr[\mathcal{M}(D) \in S]
    \leq
    e^\varepsilon \Pr[\mathcal{M}(D') \in S] + \delta .
    \label{eq:bg_dp}
\end{equation}

In DP-SGD, per-sample gradients are clipped to a norm $C$ and perturbed with
Gaussian noise:

\begin{equation}
    \tilde{g} =
    \frac{1}{B}\sum_{i=1}^{B}
    g_i \min\left(1,\frac{C}{\|g_i\|_2}\right)
    +
    \mathcal{N}(0,\sigma^2 C^2 I).
    \label{eq:bg_dpsgd}
\end{equation}

The experimental chapters show why both parameters matter. Increasing
$\sigma$ improves privacy but can damage the gradient signal; decreasing $C$
limits sensitivity and also reduces the effective noise magnitude
$\sigma \cdot C$. This interaction is central to the updated privacy benchmark
in this thesis.

\section{Root Cause Analysis and Adaptive Thresholding}
\label{sec:bg_rca_thresholding}

An anomaly detector answers whether a service looks abnormal; operators also
need to know where the problem likely started. Graph-based RCA systems build
service dependency graphs from distributed traces and rank likely causes using
structural and temporal evidence \cite{li2018microrank, zhang2019rootcause,
fu2020root}. PageRank-inspired scoring is useful because an upstream service can
explain anomalies several hops downstream \cite{page1999pagerank, brin1998anatomy}.

Thresholding remains a separate challenge. Static thresholds are simple, while
EWMA and rolling statistics adapt to local drift. Reinforcement learning frames
threshold selection as a sequential control problem in which an agent observes
recent detection outcomes and adjusts the decision boundary. Q-learning estimates
the long-term value of action $u$ in state $x$ through the Bellman update
\cite{watkins1992qlearning, sutton2018reinforcement}:

\begin{equation}
    Q(x_t,u_t) \leftarrow Q(x_t,u_t) +
    \eta \left[
    R_{t+1} + \gamma \max_{u'}Q(x_{t+1},u') - Q(x_t,u_t)
    \right].
    \label{eq:bg_qlearning}
\end{equation}

The literature motivates Q-learning for adaptive thresholding
\cite{liu2020rlthreshold}, but this thesis treats it empirically rather than
assuming success. The final evaluation shows that a tabular Q-learning tuner can
reduce final alert volume after extended interaction, but its early exploration
and reward dynamics remain unstable enough to require redesign before production
use.

The literature motivates Q-learning for threshold adaptation, but most
studies assume stationary reward signals and do not evaluate cold-start
behavior in imbalanced anomaly streams. This thesis treats Q-learning
thresholding as an empirical question (\S\ref{sec:eval_threshold}) and
reports both the eventual stabilization and the cold-start failure as
equally important outcomes.

\section{Research Gap}
\label{sec:bg_gap}

Existing work usually studies the pieces of the monitoring problem separately.
Deep anomaly detectors focus on detection accuracy, federated learning studies
distributed optimization, differential privacy studies formal leakage bounds,
and RCA systems study graph-based diagnosis. Operational monitoring, however,
requires these mechanisms to interact. A detector with high offline AUC is
insufficient if thresholds create noisy alerts, if training centralizes raw
telemetry, or if alerts do not identify a likely service-level cause.

The state of the art therefore provides strong building blocks but leaves an
integration gap. LSTM-based detectors are effective for temporal telemetry, but
many assume centralized training or offline threshold calibration. Federated
learning avoids raw-data centralization, but FL studies often report
optimization accuracy rather than alert quality, RCA usefulness, or
communication feasibility. DP-SGD provides formal protection, but its practical
utility depends on clipping, sampling, and convergence dynamics that are often
hidden behind a single $\varepsilon$ value. Graph-based RCA improves diagnosis,
but is usually evaluated separately from the detector that produces the
anomalous service set.

This thesis addresses that intersection by evaluating edge-local detection,
FedAvg, DP-SGD, compression, causal RCA, and adaptive thresholding under one
coherent architecture. Its contribution is not that each mechanism exists in
isolation, but that their combined behavior is measured as a deployable
monitoring pipeline, including the negative Q-learning cold-start result and
the updated DP clipping sensitivity benchmark.

The background surveyed here identifies the building blocks and their
individual limitations. Chapter~\ref{ch:model_architecture} narrows the focus
to the LSTM-Autoencoder, the specific detection architecture used by this
framework, and defines its structure, training procedure, and the design
choices that make it suitable for edge deployment.
